\documentclass[11pt]{article}
\usepackage[a4paper,margin=1in]{geometry}
\usepackage{amsmath,amssymb,amsthm,booktabs}
\usepackage[T1]{fontenc}
\usepackage{lmodern}
\usepackage{hyperref}

\newtheorem{theorem}{Theorem}[section]
\newtheorem{proposition}[theorem]{Proposition}
\newtheorem{corollary}[theorem]{Corollary}
\newtheorem{remark}[theorem]{Remark}
\newcommand{\Mob}{\mu}
\newcommand{\ParityFactor}{A_{\{2\}}}

\title{Parity-factor M\"obius capacity limit}
\author{RH research program}
\date{August 2026}

\begin{document}
\maketitle

\begin{abstract}
We isolate an exact constraint-graph theorem from a postcritically finite
quadratic map.  Its three-cell Markov partition has matrix
\(A=\left(\begin{smallmatrix}0&0&1\\0&0&1\\1&1&0\end{smallmatrix}\right)\),
and the binary factor \(\ParityFactor\) consists of words whose positive symbols
lie in one parity class.  For the M\"obius prefix, the resulting adaptive
capacity has an exact finite formula and satisfies
\(K_N^{(2)}/N\to4/\pi^2\).  The proof uses parity Mertens cancellation and
odd/even squarefree densities.  This is a reduced parity factor, not the
four-state distance-two constraint of RH-366.  The optimizer is chosen after
reading the arithmetic prefix, so no canonical operator, prime trace,
Riemann-zero model, Hilbert--Polya construction, or RH implication is claimed.
\end{abstract}

\section{Scope and source boundary}

The external source \cite{DynaZeta2026} proves that the restriction of
\(f_u(x)=1-ux^2\) at the real root
\[
 u_c^3-2u_c^2+2u_c-2=0,
 \qquad r=u_c-1,\qquad J=[-r,1],
\]
has the postcritical itinerary
\[
0\longmapsto1\longmapsto-r\longmapsto r\longmapsto r.
\]
The partition \(I_1=[-r,0], I_2=[0,r], I_3=[r,1]\) has transition matrix
\[
 A=\begin{pmatrix}0&0&1\\0&0&1\\1&1&0\end{pmatrix}.
 \tag{1}
\]
Labelling \(I_1\) by $1$ and the other intervals by $0$ gives the factor
\(\ParityFactor\): all $1$-positions lie in one parity class.  The same source
records the boundary-aware factor zeta function
\[
 \zeta_{f_{u_c}|_J}(z)=\zeta_{\ParityFactor}(z)=\frac{1+z}{1-2z^2},
 \tag{2}
\]
but (2) is used here only as a provenance anchor.

The distinction from RH-366 is essential.  RH-366 permits positive symbols
on both parity paths but forbids two positives at distance two; its capacity
limit is still open.  Here we study a distinct, explicitly defined reduced
factor
\(\ParityFactor\), where all positive positions must use one parity class.

\section{The finite parity-factor capacity}

For $N\geq1$, write $\Mob_n=\Mob(n)$ and
\[
 M_N=\sum_{n\leq N}\Mob_n,
 \qquad
 P_r(N)=\#\{n\leq N:n\equiv r\pmod 2,\ \Mob_n=1\},
\]
\[
 N_r(N)=\#\{n\leq N:n\equiv r\pmod 2,\ \Mob_n=-1\},
 \qquad r\in\{0,1\}.
\]
Define
\[
 K_N^{(2)}=\max_{\varepsilon\in(\ParityFactor)_N}
 \left|\sum_{n=1}^N\Mob_n\varepsilon_n\right|,
 \tag{3}
\]
where \((\ParityFactor)_N\) is the set of sign words for which the set
\(\{n:\varepsilon_n=+1\}\) is empty or contained in one residue class
modulo $2$.

\begin{proposition}[Exact finite capacity]
\label{prop:finite}
For every $N$,
\[
 K_N^{(2)}=
 \max_{r\in\{0,1\}}
 \max\bigl\{|{-M_N+2P_r(N)}|,\ |{-M_N-2N_r(N)}|\bigr\}.
 \tag{4}
\]
\end{proposition}

\begin{proof}
Start with the all-negative word, whose score is $-M_N$.  If positives are
allowed only in parity class $r$, changing a position with $\Mob_n=1$ from
$-1$ to $+1$ raises the score by $2$, while changing a position with
$\Mob_n=-1$ lowers it by $2$.  Thus the largest score in that class is
$-M_N+2P_r(N)$, obtained by flipping every positive M\"obius entry; the
smallest is $-M_N-2N_r(N)$, obtained by flipping every negative entry.  Empty
or partial flips cannot improve either extremum.  Taking absolute values and
then the two parity classes proves (4).
\end{proof}

\section{All-order limit}

We use the standard consequences of the prime number theorem in arithmetic
progressions and squarefree counting:
\[
 \frac{M_N}{N}\to0,\qquad
 \frac{P_1(N)+N_1(N)}{N}\to\frac4{\pi^2},\qquad
 \frac{P_0(N)+N_0(N)}{N}\to\frac2{\pi^2}.
 \tag{5}
\]
The signed parity cancellation is an independent input: Davenport's fixed
frequency estimate at \(\theta=1/2\) gives
\[
 A_N=\sum_{n\le N}\Mob(n)(-1)^n=o(N)
 \quad\text{(see \cite{Davenport1937}).}
\]
Together with \(M_N=o(N)\), this implies that the even and odd M\"obius
sums, \((M_N+A_N)/2\) and \((M_N-A_N)/2\), are both $o(N)$.  The
squarefree sieve gives the two densities in (5).  Since
\(P_r-N_r=o(N)\) on each class, (5) gives
\[
 \frac{P_1(N)}N,\frac{N_1(N)}N\to\frac2{\pi^2},
 \qquad
 \frac{P_0(N)}N,\frac{N_0(N)}N\to\frac1{\pi^2}.
 \tag{6}
\]

\begin{theorem}[Parity-factor capacity limit]
\label{thm:limit}
For the factor \(\ParityFactor\),
\[
 \boxed{\displaystyle \lim_{N\to\infty}\frac{K_N^{(2)}}N=\frac4{\pi^2}.}
 \tag{7}
\]
\end{theorem}

\begin{proof}
By (4) and (6), the two odd-parity candidates converge after division by
$N$ to $+4/\pi^2$ and $-4/\pi^2$, respectively.  The two even-parity
candidates converge to $+2/\pi^2$ and $-2/\pi^2$.  The maximum of their
absolute values therefore converges to $4/\pi^2$.
\end{proof}

\begin{remark}[What this does not prove]
The limit in (7) is not the unresolved RH-366 capacity for the distance-two
four-state survivor.  It is a capacity law for the distinct reduced factor
\(\ParityFactor\); the two languages are not comparable by set inclusion.
Nor does (7) turn an adaptive
sign word into a nonadaptive orbit, a canonical coupling, or a trace.
\end{remark}

\section{Frozen executable audit}

The package checks the integer M\"obius sieve, (4) against exhaustive words
through $N=12$, and the source hashes.  At $N=2^{20}$ it records
\[
 K_N^{(2)}=425095,
 \qquad K_N^{(2)}/N=0.40540218353271484.
\]
This row is a reproduction diagnostic only.  It is not a fit and is not used
to establish (7).

\section{Route and Gate ledger}

Route A is \texttt{GO}: (1)--(7) give a source-locked PCF realization, an
exact finite identity, and an all-order capacity theorem.  Route B is
\texttt{STOP\_SCOPED}: the maximizing word reads the complete M\"obius prefix,
and the quantities are scalar adaptive sums.  They are not a canonical
dynamical determinant, a von Mangoldt prime-power trace, or a completed-zeta
divisor.  Gates A--E remain false/open.  No Hilbert--Polya operator,
Riemann-zero identification, or proof of RH follows.

\begingroup\sloppy
\begin{thebibliography}{9}
\bibitem{DynaZeta2026}
Source package, \emph{A postcritically finite quadratic realization of
$W=2$}, source commit \texttt{7fd3a3fdd5a6}, 2026.
\bibitem{RH366}
RH-366, \emph{M\"obius orthogonality, adaptive encoding, and Parry
covariance}, repository release, 2026.
\bibitem{RH367}
RH-367, \emph{Boundary-aligned cyclic-Ulam structure and phase-local
leakage}, repository release, 2026.
\bibitem{Mirsky1949}
L. Mirsky, Arithmetical pattern problems relating to divisibility by
$r$th powers, \emph{Proc. London Math. Soc.} 50 (1949), 497--508.
\bibitem{Davenport1937}
H. Davenport, On some infinite series involving arithmetical functions (II),
\emph{Quart. J. Math.} 8 (1937), 313--320.
\end{thebibliography}
\endgroup

\end{document}
